% !PRINTED 649
\weberpage{649}
%<<CONT

If two holomorphic functions
\[
\begin{aligned}
\alpha &= \alpha_0 x^h + \alpha_1 x^{h-1} + \cdots + \alpha_h, \\
\beta &= \beta_0 x^k + \beta_1 x^{k-1} + \cdots + \beta_k
\end{aligned}
\]
have a product
\[
\gamma = \gamma_0 x^{h+k} + \gamma_1 x^{h+k-1} + \cdots + \gamma_{h+k},
\]
whose coefficients are integral functionals, then the $\gamma_0$,
$\gamma_1, \ldots, \gamma_{h+k}$ can all be divisible by a prime functional $\pi$
only if either all of $\alpha_0, \alpha_1, \ldots, \alpha_h$, or all of $\beta_0, \beta_1, \ldots, \beta_k$
are divisible by $\pi$, which is proved in exactly the same way as in Bd. I, § 2.

If now (2) is an integral functional without $\varphi_0, \varphi_1, \ldots, \varphi_m$
being integral, then one can determine an integral functional $\mu$, and in it a
prime factor $\pi$, in such a way that
\weq{3}{\chi = \mu\varphi = \alpha_0 x^m + \alpha_1 x^{m-1} + \cdots + \alpha_m,}
is divisible by $\pi$, while $\alpha_0, \alpha_1, \ldots, \alpha_n$ are not all divisible by $\pi$.
% [sic] printed $\alpha_n$, the sum in (3) runs to $\alpha_m$

If now $\varphi$ is integral, then it satisfies an equation:
\weq{4}{E\varphi^m = a_1 E_1 \varphi^{m-1} + a_2 E_2 \varphi^{m-2} + \cdots + a_m E_m}
and from this, by multiplication by $\mu^m$:
\weq{5}{E\chi^m = \mu(a_1 E_1 \chi^{m-1} + a_2 E_2 \mu \chi^{m-2} + a_m E_m),}
% [sic] printed $\mu(a_1 E_1 \chi^{m-1} + a_2 E_2 \mu \chi^{m-2} + a_m E_m)$, multiplying (4) by $\mu^m$ gives $\mu(a_1 E_1 \chi^{m-1} + a_2 E_2 \mu \chi^{m-2} + \cdots + a_m E_m \mu^{m-1})$; the $+ \cdots$ and the factor $\mu^{m-1}$ of the last term are missing
in which the $a_1, a_2, \ldots$ are integral functions in $Z$, and the $E, E_1, \ldots, E_m$
are units in $\bar{Z}$. The units $E, E_1, \ldots, E_m$ may indeed still
contain $x$, but they may be assumed to be holomorphic.
If we arrange the right-hand and the left-hand side of (5) according to $x$, then
the coefficients of $E$ are not all divisible by $\pi$. Likewise,
by the theorem mentioned, the coefficients of $\chi^m$ are not
all divisible by $\pi$, while on the right-hand side all coefficients
have the factor $\mu$, and hence also the factor $\pi$. That is
impossible, and with this our theorem is proved (Bd. II, § 159).

8. \gesperrt{A holomorphic integral functional $\varphi$ is the
greatest common divisor of all its coefficients.}

If $\alpha, \beta, \ldots$ are the coefficients of $\varphi$, then $\varphi$ is in any case
divisible by the greatest common divisor of $\alpha, \beta, \ldots$.
If we decompose a linear function $z - c$ into its functional prime
factors, then in these prime functionals $\pi$ we may denote the variables
as we please, and since every prime functional must divide some $z - c$,
we may assume that the variables in it are different from the
variables $x, y, \ldots$ in $\varphi$. Accordingly we can
%CONT>>
